% !TeX root = ../main.tex

\chapter{Data Ethics and Governance}
\label{apx:ethics}

Every result in this dissertation is measured on check-ins published by location-based
social networks, and neither collection used here was assembled by the author. This
appendix states where each came from, under what terms, what the pipeline does with it,
and what remains unsettled.

\section{Where the data came from}

The Gowalla check-ins are read from a category-annotated deposit published on Figshare,
DOI \path{10.6084/m9.figshare.22126586.v2}, which carries the Creative Commons CC0
public-domain dedication.\footnote{\url{https://creativecommons.org/publicdomain/zero/1.0/}}
Four of its files enter the pipeline: the check-in table, two place tables, and the
category structure file. The seven place categories used throughout the dissertation are
the top-level names in that structure file, so the taxonomy arrives with the deposit;
what the pipeline adds is the mapping of each finer category up to its top-level parent,
plus a short local supplement for names the structure file leaves uncovered.

Two qualifications belong with that label, and both come from the record itself. The
dedication was applied by the depositor of this copy, identified there by a single name,
and not by the party that collected the data. The record also names a further site as
the origin of the files, and that address now redirects to an unrelated commercial site,
so its terms could not be read. What is supportable is therefore a statement about the copy in hand,
which is distributed under CC0; that Gowalla data as such is in the public domain is a
broader claim, and nothing that could be opened supports it. A second and older Gowalla
release, hosted by the Stanford Large Network Dataset Collection and cited here as the
collection reference \cite{cho2011gowalla,jure2014snap}, is a different artifact: it
carries no place categories, and its page states no license.

The Istanbul check-ins come from the Massive-STEPS collection
\cite{wongso2025massivesteps}, distributed on Hugging Face under the Apache License
2.0.\footnote{\url{https://huggingface.co/datasets/CRUISEResearchGroup/Massive-STEPS-Istanbul}}
The project repository carries the same license. Massive-STEPS is derived material, and
its documentation names the Semantic Trails Dataset and Foursquare Open Source Places as
its sources; among the records that could be opened, no upstream term more restrictive
than the Apache License 2.0 appeared. Two facts qualify that. The Foursquare
distribution is additionally access-gated behind an agreement prompt, so a reader who goes
to that upstream must accept terms before downloading, while the Massive-STEPS copy used
here is not gated. And one check is outstanding: the Foursquare product terms themselves
were not read, only the license tag on the distribution.

\section{Real people, and how the traces are handled}

A user identifier in these files is a pseudonym and not a name, but a sequence of
timestamped visits is still a description of one person's movements. Pseudonymity is not
anonymity. The mobility literature treats the residual risk as open, and reports that
models trained on mobility data raise privacy questions during training as well as at
prediction time, because the portions of information a model absorbs cannot be
controlled directly \cite{luca2021mobilitysurvey}. Public availability settles who may
hold the files, not what can be inferred from them.

Against that, this work adds no de-identification of its own. No coordinate is
perturbed, rounded, generalized, or masked, and no formal privacy mechanism is
applied. Latitude and
longitude are used at the precision the sources publish, because the prediction target
of Chapter~\ref{ch:mobiwac} is itself spatial, and a coarsened coordinate would change
the quantity being measured.

Exposure is limited by restriction instead. The Figshare deposit also contains a
social-graph file and a user-profile file; neither is declared as an input, and no code
reads either one, so the friendship network and the profile fields never enter the
study. User identifiers are carried as opaque integers exactly as the sources supply
them, and a non-numeric identifier is replaced by a position index that is not kept.
Nothing in the code links a user across the two collections or to any outside source.

No check-in data is redistributed. The public repository publishes code and
configuration, and the data directory is excluded from version control, so a reader
reproducing the results obtains both collections from the sources named above.
